% Français 9115, batch 30 — Gallica views 581–600.
% Pass: Opus 5.5 (claude-opus-5-5), 2026-09-22 — first pass, unchecked against the views by a human.
%
% All twenty views are on disk in the mirror; none is missing.
%
% What the twenty views hold. Four pieces, all on the vibrating elastic
% surfaces of the Institut's prize question, none of them number theory.
%
%   v. 581–592  The end of a long memoir, begun before this batch: her
%               pages (154) to (165), written recto and verso, a fair hand
%               with corrections between the lines. It continues v. 580,
%               whose page (153) ends « comme dans la bande » (seen on the
%               thumbnail only, not transcribed); v. 581 opens « droite, et
%               lorsqu'[elle, struck] la courbure est uniforme ». Subject: the ring and
%               the curved lamina against Euler's theory and Chladni's
%               experiments — Chladni's metal band (« traité d'acoustique »
%               N.o 89), Euler's « N.o 54 de son M. sur la lame », which
%               excludes all but supported ends for the ring; then N.o 25
%               (v. 581), her own experiments with laminae cut from glass
%               (forks, rings of 83 mm diameter and widths 9, 20 and 35 mm,
%               a circular plate and a ring differing by an 11 mm hole),
%               with the notes heard (si 1, mi 2, fa# 3 …) and their
%               intervals against Euler's numbers; the conclusion that
%               curvature does not alter the modes, only a varying
%               elasticity does, and that Euler wrongly excluded fixed ends
%               for the ring; the equation z = … A sin(beta phi + B), that of the
%               tensed string, extended to lamina, ring, « tambour
%               circulaire », circular plate and hemispherical bell, with
%               Chladni's tables (p. 185?, p. 236). The memoir ends on v.
%               592, a third of the way down page (165): « Je termine enfin
%               ce long mémoire … les conditions du programme proposé par
%               la classe. » Taken with the brief (prize memoir at page
%               (130), N.o 20, on v. 540) this is presumably the end of
%               that memoir; the pages between were not read in this pass.
%   v. 593–595  A fresh start, in a larger, very regular hand: the title
%               « Mémoire sur la question proposée par la première classe
%               de l'institut », the question quoted, the Bacon epigraph
%               (« Baco Novum organum, lib. 1, n. XCII. tom. 1 p. 298,
%               edit. in f.o London 1740 »), and the opening of the memoir
%               — a fair copy of the draft read at v. 398–399, still with
%               old lines hatched out between the new ones — then, on the
%               back (page (2), v. 595), « §. Recherche de l'équation
%               différentielle dela surface élastique vibrante » and N.o 1,
%               stopping mid-sentence on « le cosinus de l'angle ». View
%               594 is v. 593 photographed a second time (same text, same
%               hatching, same ink 292; only the framing differs) and is not
%               transcribed.
%   v. 596–597  Pages (23)–(24) in the same regular hand, one leaf recto
%               and verso: the rectangular plate with its sides parallel to
%               the axes; the boundary conditions of the sides and the
%               « quatre couples d'équations »; « points de double limite »
%               at the corners. Begins mid-sentence (« équations,
%               j'examinerai ») and ends mid-sentence (« on ait dy = 0, »);
%               neither neighbour is in this batch. Whether (23)–(24) belong
%               with the (1)–(2) of v. 593–595 is suggested by the hand and
%               the subject, not settled.
%   v. 598      Another state of the same opening, unpaginated and heavily
%               corrected — « Mémoire sur cette question … », « Si on
%               vouloit embrasser d'abord cette question dans toute sa
%               généralité », « une équation différentielle du quatrième
%               degré » (the last words struck), a long margin addition —
%               breaking off two-thirds down the leaf on « en substituant
%               aux. ». Nothing orders it against v. 593 and v. 398.
%   v. 599      The blank back of v. 598 (pencil strokes on the mount only);
%               not transcribed.
%   v. 600      Page (61), struck, of yet another run: the integrals
%               labelled with small crosses (« (++) »), figures of Chladni
%               (« fig. 796? et 86 », m = 5, m = 6), a long oblique stroke
%               cancelling two paragraphs, a margin note running on across
%               the foot of the page. Its first line (« En effet j'ai vu
%               bien des fois ») continues a page not in this batch; its
%               last word, « forme », runs on to v. 601, whose thumbnail
%               shows page (62) (« En effet, si on fait b = 0, n = 1 et
%               m = 5, l'intégrale (…) deviendra »).
%
% The numbers on the leaves, and why no \folio{} is written. (1) The
% top-right ink series, one number per physical leaf, on the recto — here
% the side her pagination makes even, (154), (156) … — or on the only
% written side: 286 (v. 581), 287
% (583), 288 (585), 289 (587), 290 (589), 291 (591), 292 (593, and again
% on the repeat 594), 293 (596), 294 (598), 295 (600) — no gap, the
% written versos 582, 584 … 592, 595, 597 carrying none. Penned, as
% everywhere in this volume, so no \folio{} is written (hand.md, « The open
% question of this volume »). Laubenbacher and Pengelley's ff. 348r–349r
% are not reached: the series stands at 295 on v. 600. (2) Her page
% numbers, top centre, in parentheses: (154)–(165) on v. 581–592, the pages
% on the rectos, (154), (156) … (164), struck obliquely and those on the
% versos not — the alternation of v. 429–439 with the parity reversed; none
% on v. 593 and 598; (2) on v. 595; (23), (24) on v. 596–597; (61), struck,
% on v. 600, the second figure blotted (v. 601's (62) supports it). (3) A
% paragraph number, N.o 25, inside the line at v. 581. (4) Numbers cited:
% « N.o 89 du traité d'acoustique » (v. 581), Euler's « N.o 54 » (v. 583),
% her own « N.os 22 et 23 » (v. 587), Chladni's « table p. 185 » (last
% figure doubtful) and « table p. 236 » (v. 591), « fig. 796 et 86 »
% (v. 600, the first doubtful). Every number was read on its view; none is
% inferred.
%
% Notation and conventions. Her spelling is kept and not flagged: « égals
% », « satisfesante », « crû », « appuyés » for the ends, « fixés »,
% « fesant », « rejetter », « hazard », « quarrés », « je prens »,
% « arreterai », « dela », « ceque », « cequi », « aunombre », « Il ya »,
% « parceque », « parconséquent ». Notes are written as she writes them,
% name then octave figure: « si 1 », « mi. 2 », $\text{fa}^{\sharp}$ 3.
% Her double integral « SS » and contour sum « S » are kept as S; the
% variation sign is $\delta$; her braces $\{\ \}$ are kept; « ddz » is
% kept. Letters underlined in formulas stay \underline in maths;
% underlined words in running text are given as \emph{} (text-mode
% \underline fails the TEI check). The capital pi of the integrals is set
% $\Pi$, as batches 20–22 did. Plurals cancelled with a small cross on the
% final letter are given either in the singular with a note or with the
% letter in \struck{} (l'anneau\struck{x}, diamètre\struck{x}).
%
% What could not be read. \ill{} stands in three places: a struck figure
% after « 11 » (v. 589); at v. 600, a struck word before « se réduisent »
% and the first word of a struck group in the margin note. Readings
% offered with doubt (\uncertain{}): « La fi » and the octave figure « 3 »
% (v. 585); « 35 » millimetres (v. 588); « 185 » (v. 591); « nombre » and
% « d'essais » under hatching (v. 593); « fort » and « habile » (v. 598);
% « 796 » and the second integral's label (v. 600). The order of the last
% words of v. 600's margin note is offered, not assured.
%
% Nothing here was checked against any published edition or study.
\documentclass[11pt,a4paper]{article}
% The preamble shared by every transcription of the mathematicians' manuscripts in Gallica.
%
% It rests on one idea: **uncertainty compounds, it does not resolve.** A
% transcription that smooths over a doubtful word has destroyed the only thing
% it was made for — a reader must be able to tell, at every point, what was on
% the page and what was supplied. The apparatus macros below are therefore the
% critical apparatus, not decoration.
%
% The same commands are recognised by `scripts/render.mjs`, which produces the
% site's reading view, and by `scripts/tei.mjs`. A macro added here must be
% added there too, or rendering fails loudly — which is the wanted behaviour.
%
% Comments here are English, like the rest of the repository. The documents
% this preamble sets are French, and that is deliberate: see the skills.

\ProvidesPackage{germain}[2026/09/15 Transcriptions of mathematicians' manuscripts in Gallica]

\usepackage{amsmath,amssymb,amsthm}
\usepackage{mathrsfs}
% Kept from the parent preamble so the renderer's diagram support stays
% available; these manuscripts draw no commutative diagrams, and nothing here uses it.
\usepackage{tikz-cd}
\usepackage[english,french]{babel}
\usepackage{fontspec}

% --- Greek ------------------------------------------------------------------
% The text face stays fontspec's default, Latin Modern, which has no Greek:
% XeTeX drops every glyph it lacks with a line in the log and nothing on the
% page, and the Greek words the manuscripts quote vanished from the PDFs. So
% the Greek and Greek Extended blocks get a XeTeX character class of their own,
% and entering or leaving a run of it switches the family to CMU Serif — the
% same Computer Modern design, with polytonic Greek — and back. Only the
% family changes: a Greek word in \textit or \textbf keeps its shape and series.
% The transitions to and from every other class carry the tokens that class
% has towards an ordinary letter, so French spacing before « ; » or inside
% « » still holds after a Greek word. No grouping, so no brace or \endgroup
% can come out unbalanced; maths is left alone, since XeTeX inserts nothing
% there. Kept identical in `exercices.sty`.
\makeatletter
\newfontfamily\germainGreekfont{cmun}[
  Extension=.otf, NFSSFamily=germaingreek,
  UprightFont=*rm, ItalicFont=*ti, SlantedFont=*sl,
  BoldFont=*bx, BoldItalicFont=*bi, BoldSlantedFont=*bl]
\newXeTeXintercharclass\germain@greek
\def\germain@greekrange#1#2{%
  \@tempcnta=#1\relax
  \loop
    \XeTeXcharclass\@tempcnta=\germain@greek
    \ifnum\@tempcnta<#2\advance\@tempcnta\@ne\repeat}
\germain@greekrange{"0370}{"03FF}
\germain@greekrange{"1F00}{"1FFF}
\def\germain@greekprev{\rmdefault}
\def\germain@greekin{%
  \edef\germain@greekprev{\f@family}\fontfamily{germaingreek}\selectfont}
\def\germain@greekout{\fontfamily{\germain@greekprev}\selectfont}
\def\germain@greekjoin{%
  \XeTeXinterchartokenstate=\@ne
  \@tempcnta=\z@
  \loop
    \ifnum\@tempcnta=\germain@greek\else
      \XeTeXinterchartoks\germain@greek\@tempcnta=\expandafter{%
        \expandafter\germain@greekout\the\XeTeXinterchartoks\z@\@tempcnta}%
      \XeTeXinterchartoks\@tempcnta\germain@greek=\expandafter{%
        \the\XeTeXinterchartoks\@tempcnta\z@\germain@greekin}%
    \fi
    \ifnum\@tempcnta<4095\advance\@tempcnta\@ne\repeat}
% babel-french sets its own classes, and saves and restores the interchar
% state, each time a language is selected: join again after it does.
\addto\extrasfrench{\germain@greekjoin}
\addto\extrasenglish{\germain@greekjoin}
\AtBeginDocument{\germain@greekjoin}
\makeatother

\usepackage{geometry}
\geometry{a4paper,margin=25mm,marginparwidth=28mm}
\usepackage{marginnote}
\usepackage{xcolor}
\usepackage[normalem]{ulem}
\setlength{\emergencystretch}{3em}
\usepackage{hyperref}
\hypersetup{colorlinks=true,linkcolor=black,urlcolor=black,pdfborder={0 0 0}}

% --- Batch metadata ---------------------------------------------------------
% Taken as they stand by the rendering script, which shows them at the head of
% the reading view. They fix what the file claims to cover. The macro names are
% the parent project's, so that the scripts stay shared; \folder here means the
% volume's slug — `fr-9115` — and \pages counts Gallica views.

\newcommand{\folder}[1]{\gdef\grothFolder{#1}}
\newcommand{\batch}[1]{\gdef\grothBatch{#1}}
\newcommand{\pages}[2]{\gdef\grothFirst{#1}\gdef\grothLast{#2}}
\newcommand{\foldertitle}[1]{\gdef\grothFoldertitle{#1}}

% The holder's own shelfmark, printed as they write it: « Français 9115 ».
\newcommand{\shelfmark}[1]{\gdef\grothShelfmark{#1}}

% The dating, copied verbatim from the holder's catalogue — never inferred
% here. « XIXe siècle » is the BnF's; a pass that dates a leaf from its content
% says so in a \note{}, as its own inference.
\newcommand{\dating}[1]{\gdef\grothDating{#1}}

% The status notice, printed in the title block and repeated as a diagonal
% watermark on every page of the PDF. The manuscripts are in the public
% domain; what this line declares is not a right but a quality — an
% unverified first pass by a machine. The skills write the line; a file
% without it gets no watermark, so leaving it out is visible in the output.
\newcommand{\watermark}[1]{\gdef\grothWatermark{#1}}

\gdef\grothFolder{?}\gdef\grothBatch{}\gdef\grothFirst{?}\gdef\grothLast{?}%
\gdef\grothFoldertitle{}\gdef\grothShelfmark{}\gdef\grothDating{}\gdef\grothWatermark{}

% --- The résumé -------------------------------------------------------------
\newenvironment{resume}
  {\small\setlength{\parskip}{0.45em}}
  {\par\medskip\hrule\medskip}

% --- Keywords ---------------------------------------------------------------
% In English, closing the modernised reading's résumé: the modern vocabulary
% under which to search for what the leaves construct. The manifest extracts
% them to tag the volume — this line is the single source of the tags.
\newcommand{\keywords}[1]{\par\smallskip\noindent
  {\footnotesize\textsc{Keywords} --- \textit{#1}}\par}

% --- The critical apparatus -------------------------------------------------

% The Gallica view number, in the margin. This is the anchor that lets a
% disagreement over a formula be settled by looking at *one* image rather than
% a volume of 750. The site uses it to turn the facsimile as the transcription
% is scrolled. A view is one scanned image: usually one side of a leaf.
\newcommand{\page}[1]{%
  \marginnote{\footnotesize\color{gray!70}\textsf{v.\,#1}}%
  \label{page:#1}\ignorespaces}

% The BnF's pencilled foliation, where the leaf carries it and a pass can read
% it: « 198r ». This is what the literature cites — « ff. 198r–208v » — and
% the one line that turns a citation into a view. Recorded, never inferred:
% a folio a pass cannot read is not written.
\newcommand{\folio}[1]{%
  \marginnote{\footnotesize\color{gray!50}\textsf{f.\,#1}}\ignorespaces}

% The modernised reading's coarser counterpart to \page{}: which views a
% section is drawn from.
\newcommand{\pagerange}[2]{%
  \marginnote{\footnotesize\color{gray!70}\textsf{v.\,#1--#2}}%
  \ignorespaces}

% Illegible: never guessed.
\newcommand{\ill}{\textcolor{red!60!black}{[\,\ldots\,]}}

% A doubtful reading: the word is given, and flagged as uncertain.
\newcommand{\uncertain}[1]{\uline{#1}}

% An editorial addition — a missing letter, an implied word.
\newcommand{\add}[1]{\textcolor{blue!50!black}{[#1]}}

% Struck out by the writer. What was crossed out is part of the document.
\newcommand{\struck}[1]{\sout{#1}}

% The transcriber's note, on the state of the page or the mathematics.
\newcommand{\note}[1]{\par\smallskip{\small\color{gray!60!black}\textit{[#1]}}\par\smallskip}

% A marginal note of hers, distinct from the body of the text.
\newcommand{\marginal}[1]{\par\smallskip\noindent\hspace{1em}%
  {\small\color{gray!60!black}#1}\par\smallskip}

% --- Title ------------------------------------------------------------------
% The title block is French, like the documents themselves.

\AddToHook{shipout/background}{%
  \ifx\grothWatermark\empty\else
    \begin{tikzpicture}[remember picture, overlay]
      \node[rotate=52, scale=3.4, align=center, text=gray!18,
            font=\sffamily\bfseries]
        at (current page.center) {\grothWatermark};
    \end{tikzpicture}%
  \fi}

\AtBeginDocument{%
  \begin{center}
    {\large\bfseries Manuscrits de mathématiciens (Gallica) ---
      \ifx\grothShelfmark\empty\grothFolder\else\grothShelfmark\fi}\\[2pt]
    {\itshape\grothFoldertitle}\\[4pt]
    {\small vues \grothFirst--\grothLast
      \ifx\grothBatch\empty\else\ (lot \grothBatch)\fi}\\[2pt]
    \ifx\grothDating\empty\else
      {\small\color{gray!60!black}Datation du catalogue : \grothDating}\\[2pt]
    \fi
    \ifx\grothWatermark\empty\else
      {\small\bfseries\color{red!45!gray}\grothWatermark}\\[2pt]
    \fi
    {\footnotesize\color{gray!60!black}Transcription établie d'après les images IIIF de Gallica.
     Source gallica.bnf.fr / Bibliothèque nationale de France.\\
     Les lectures incertaines sont soulignées, les illisibles notées [\,\ldots\,],
     les ajouts entre crochets ; \textsf{v.} numérote les vues de Gallica,
     \textsf{f.} la foliotation au crayon de la BnF.}
  \end{center}
  \vspace{1em}}

\endinput


\folder{fr-9115}
\shelfmark{Français 9115}
\batch{30}
\pages{581}{600}
\dating{1801-1900}
\watermark{Lecture automatique\\non vérifiée}
\foldertitle{Papiers de Sophie GERMAIN. — Recueil de dissertations et problèmes mathématiques et physiques.}

\begin{document}

\note{Numérotation des feuillets. En haut à droite, à l'encre, une série d'un nombre par feuillet : 286 (vue 581), 287 (583), 288 (585), 289 (587), 290 (589), 291 (591), 292 (593 et sa reprise 594), 293 (596), 294 (598), 295 (600), sans lacune ; les revers écrits n'en portent pas. Elle occupe la place de la foliotation de la Bibliothèque mais est tracée à la plume : aucun « f. » n'est donc porté. En tête, au milieu, la pagination de l'auteur, entre parenthèses : (154) à (165) aux vues 581–592, les pages des rectos, (154), (156) … (164), barrées ; (2) à la vue 595 ; (23) et (24) aux vues 596–597 ; (61), barré, à la vue 600. La vue 594 reprend la vue 593 et la vue 599 est le revers blanc du feuillet de la vue 598 : elles ne sont pas transcrites.}

\page{581}

\note{En tête, au milieu, entre parenthèses et barré : (154). En haut à droite, à l'encre : 286.}

droite, et lorsqu'\struck{elle} la courbure est uniforme, comme dans
la bande ployé à la forme d'anneau circulaire. Mais
l'élasticité, quoique constante dans tous les points de l'anneau, doit
être diminuée par l'effet de la courbure. Ainsi l'anneau doit se \struck{partag}
partager en un certain nombre de parties égals entre elles ; mais
les sons correspondans doivent être plus graves que ceux dela
bande ramené à la forme rectiligne. La différence des sons donnés
par l'anneau et par la bande droite, doit être proportionnelle
à l'amplitude de l'arc \struck{séparé} \add{terminé} par deux points de repos
consécutifs, parceque la courbure est d'autant moindre qu'il
s'agit d'un arc plus petit.

Rien n'est encore plus conforme aux phénomènes rapportés
N.o 89 du traité d'acoustique ; car M.r Chladni a observé
que les longueurs des parties vibrantes sont égales entre elles ;
et les nombres auxquels il a trouvé que les sons doivent
être proportionnels sont à ceux indiqués par la théorie
d'Euler $::2\beta-1:2\beta$. On voit donc que la différence, savoir
$\frac{(2\beta)^2}{(2\beta-1)^2}$, finit par être insensible, lorsque $\beta$ est très grand.
\note{« ramené » : le mot se lit « ramimé » ; lecture par le sens. « entre elles » est écrit d'un trait, le r et le e liés. « d'Euler » : le nom est écrit de son E en forme de L, comme ailleurs dans le volume. La fraction est écrite comme ci-dessus ; c'est un rapport qu'elle nomme « la différence ».}

N.o 25. Quelque satisfesante que me paroisse la manière dont j'ai
crû pouvoir concilier la théorie d'Euler et le résultat des
expériences de M.r Chladni ; l'importance d'une question dans laquelle
il ne s'agit de rien moins que de décider si la courbure du
corps sonore permet \struck{encore} de lui appliquer \add{encore} l'analyse qui
\note{Le Q de « Quelque » est écrit par-dessus une autre majuscule.}

\page{582}

\note{En tête, au milieu, entre parenthèses : (155), le dernier chiffre repassé ; ce nombre n'est pas barré. Aucun nombre en haut à droite : c'est le verso écrit du feuillet 286.}

convient au cas où le même corps est supposé rectiligne,
m'a engagé à consulter moi même l'expérience, mais en
écartant la cause que je supposois avoir influé sur les
faits observés par M.r Chladni.

Avant d'entrer en matière, il faut observer que, si les
effets de la courbure sont nuls, le cas de la courbe rentrante
ne peut différer de ceux \struck{dans} lesquels la lame affecte
toute autre forme, qu'en ceque les deux extrémités réunies
alors dans un seul et même point, sont nécessairement
dans un état semblable.

On pourra donc regarder le point d'application des
doigts sur la lame élastique circulaire, tantot comme
un point fixé où seront réunies les deux extrémités
égalment fixés, tantot comme un point d'appui où seront
réunies les deux extrémités appuyés. Rien n'empêchera
même, en cessant de prendre l'origine au point d'application,
de considerer ce point comme un simple point de
repos appartenant au cas des deux extrémités libres,
extrémités qui pourront être censés réunis dans un
des points mobiles de l'anneau.
\note{« pourra » : les deux dernières lettres sont écrites par-dessus d'autres, et « doigts » a son s final repassé. « réunies » (troisième ligne de l'alinéa) a ses dernières lettres repassées.}

Je ne puis dissimuler pourtant qu'Euler paroît être
à l'égard des effets de la courbure considérés dans
l'anneau, d'une opinion toute contraire à celle que
je viens d'exposer ; car ce grand géomètre, en

\page{583}

\note{En tête, au milieu, entre parenthèses et barré : (156). En haut à droite, à l'encre : 287. Le feuillet est monté plus bas sur sa feuille que le précédent.}

prenant successivement $s=a$, $s=2a$ ; $s=3a$ \&c, arrive à conclure
$\alpha=0$, $\beta=0$ et $\delta=0$, c'est à dire qu'il exclut tout autre cas
que celui des extrémités appuyés (V. N.o 54 de son M. sur la lame)
\note{La première lettre de « $\alpha=0$ » a la forme d'un « u », comme l'$\alpha$ de la copie latine d'Euler. « M. » : abrégé, sans doute pour « mémoire ». Le 4 de « 54 » a la forme d'un « b ».}

Malgré le respect dû à l'autorité \struck{de si illustre auteur} \struck{d'un auteur aussi} \add{de} \struck{cet illustre auteur} \add{ce savant auteur}, je
ne puis m'empêcher d'observer que les suppositions $s=a$, $s=2a$ \&c.
doivent être regardées comme purement mathématiques ; qu'excepté
la première, elles n'ont aucune existence réelle et physique,
et qu'elles sont parconséquent étrangères à la question.
\note{Trois leçons se superposent au-dessus de « l'autorité » : sur la ligne, « de si illustre auteur », barré ; au-dessus, « d'un auteur aussi », barré ; plus haut, « de cet illustre auteur », dont « cet illustre auteur » est barré ; enfin, tout en haut, « ce savant auteur », non barré. « réelle » et « physique » portent chacun une petite croix sur la lettre finale : le pluriel annulé, qu'on donne ici au singulier.}

S'il en étoit autrement, il faudroit qu'un anneau tel petit qu'on
voulut le supposer, pût être équivalent à une lame droite
d'une longueur infinie, c'est à dire qu'il faudroit que l'intensité
du son, qui, toutes choses égales d'ailleurs, dépend du nombre
des parties vibrantes à la fois, fut égal de part et d'autre.
Je reviens aux expériences que j'ai tentées.
\note{Une petite croix est tracée sur l's final de « égales ».}

Dans la vue d'écarter l'influence physique de la courbure,
et d'obtenir une élasticité constante dans tous les points de
la lame ; au lieu d'employer, comme l'a fait M.r Chladni,
une bande métallique susceptible de prendre différentes formes,
j'ai fait tailler dans un morceau de verre bien choisi
des lames de diverses courbures.

J'ai d'abord essayé les lames composées de deux branches
droites réunies par une partie circulaire, j'ai reconnu que,
dans ce cas, qui est celui des fourches, les points de repos
\note{Une petite croix est tracée sous « que », à la fin de l'avant-dernière ligne ; aucune note ne lui répond sur cette page.}

\page{584}

\note{En tête, au milieu, entre parenthèses : (157), non barré. Aucun nombre en haut à droite : c'est le verso écrit du feuillet 287.}

étoient d'autant plus distants les uns des autres, qu'ils étoient plus
voisins du milieu dela lame qui est en même tems le milieu
de la partie courbe de cette lame. La mesure doit être
prise sur les parties extérieures dela courbe ; car, la largeur
de la lame n'influant pas sur le son, il depend uniquement
de l'étendue de son contour extérieur : Cette observation
s'applique à tous les cas suivans ; ainsi elle sera toujours
sous entendue.

Le résultat que j'ai obtenu est, comme on voit, entièrement
contraire à celui qu'ont donné les expériences de M.r Chladni.
les miennes font voir que la fourche d'une élasticité constante,
vibre absolument de la même manière qu'une lame droite,
et que parconséquent c'est le cas des extrémités libres
en vertu des conditions $\frac{ddz}{dx^2}=0$ et $\frac{d^3z}{dx^3}=0$ qui s'y
montrent le plus aisément.
\note{« donné » : la finale est repassée, peut-être un « s » effacé.}

Au contraire, lorsque l'élasticité n'est pas la même dans
tous les points dela lame, non seulement la distance
entre les points d'appui varie en raison de ce coëfficient,
mais encore, suivant l'observation de M.r Chladni,
cette fourche vibre alors comme une lame dont les deux
extrémités seroient appuyés, ou, cequi revient au même,
comme une lame dont les extrémités seroient libres
en vertu des conditions $\frac{dz}{dx}=0$ et $\frac{d^3z}{dx^3}=0$.
\note{La première condition est écrite ici $\frac{dz}{dx}=0$, sans exposant, là où l'alinéa précédent écrivait $\frac{ddz}{dx^2}=0$ ; on la donne telle qu'elle est écrite. Le bas de la page est blanc.}

\page{585}

\note{En tête, au milieu, entre parenthèses et barré : (158). En haut à droite, à l'encre : 288.}

Peut-être ce cas ne se montre-t-il sur la fourche \struck{qu'à cause du}
soumise à l'expérience par M.r Chladni, qu'à cause du grand nombre
de points de limites auxquels il donne lieu, et qui, en mettant
chaque partie qu'ils séparent dans la même position analytique
que si elles étoient réellement isolées, permet de considérer
l'élasticité comme uniforme dans chacune d'elles.
\note{« partie » : une lettre finale, sans doute un s, est barrée d'une petite croix.}

A l'égard du cas de l'anneau, j'ai fait également tailler
des lames de verre de forme circulaire et de grandeurs
variée. La première que j'ai essayé, étoit large de
9 millimètres environ, et son diamètre étoit de 83 millimètres.
J'ai aisement obtenu trois figures différentes : la première
présentoit trois points de repos inégalement espacés, et
aunombre desquels étoit compris le point d'application
des doigts : Le son qui accompagnoit cette figure
étoit si.~1.. La seconde figure étoit composée de quatre
points de repos situés à égale distance, les uns des autres.
le son qui l'accompagnoit étoit mi.~2. \uncertain{La fi}
la troisième figure étoit formée de 6 points de repos
pareillement situés à égale distance, les uns des autres, et
le son correspondant étoit $\text{fa}^{\sharp}$~\uncertain{3}.
\note{Les sons sont notés par le nom de la note suivi d'un chiffre, qui paraît désigner l'octave ; le dièse est écrit au-dessus de « fa ». Le dernier chiffre a la forme d'un 9 à queue courte ; la suite 1, 2 fait offrir 3, avec doute. En fin de ligne, après « mi. 2. », deux mots commencés, lus « La fi », non barrés, que la ligne suivante reprend par « la troisième figure ».}

Il est aisé de voir que les deux dernières figures appartiennent
au cas des extrémités appuyés ; car on les explique en fesant

\page{586}

\note{En tête, au milieu, entre parenthèses : (159), non barré. Aucun nombre en haut à droite : c'est le verso écrit du feuillet 288.}

pour l'une $n=4$ pour l'autre $n=6$ ; et les sons qui doivent être
alors entr'eux $::16:36$ donnent effectivement $\frac{9}{4}$, c'est adire,
1 octave et 1 ton pour leur intervalle, précisément comme l'indique
l'expérience.
\note{« donnent » : la finale est écrite par-dessus une autre.}

A l'égard dela première figure, puisqu'elle présente un nombre
impair de points de repos, et que ces points sont inégalement
distans les uns des autres, elle ne peut appartenir au cas
des extrémités appuyés. Le son qui l'accompagne annonce
qu'elle doit être rapportée au cas des extrémités fixés.
En effet, d'après la théorie d'Euler, le nombre auquel le
son doit être proportionnel lorsqu'il y a deux points
de repos dans la lame droite dont les deux
extrémités sont \struck{appuyés} \add{fixés}, est $\frac{49}{4}:\frac{64}{49}$ doit donc exprimer
l'intervalle entre les sons qui conviennent à ce cas
et à celui de 3 points de repos dans la lame droite
dont les extrémités sont appuyés ou de 4 points de
repos dans l'anneau élastique ; or $\frac{64}{49}$ diffère à
peine de $\frac{4}{3}$, c'est adire de 2 tons et $\frac{1}{2}$ ton
qui est en effet l'intervalle que donne l'expérience
puisqu'un de ces sons étant si~1, l'autre \struck{est}
\struck{trouvé} \struck{est} mi.~2.
\note{« est » est écrit « ett », comme ailleurs sur ces feuillets. Entre $\frac{49}{4}$ et $\frac{64}{49}$ le signe est un deux-points traversé d'un petit trait, suivi d'un point-virgule ; on le donne par un deux-points. Une petite croix suit $\frac{64}{49}$ ; aucune note ne lui répond sur la page. La phrase semble avoir perdu un mot entre « $\frac{49}{4}$ » et « $\frac{64}{49}$ doit donc » ; elle est donnée telle qu'elle est écrite.}

Le petit nombre d'expériences que je viens de rapporter
suffit pour prouver que, conformément à la doctrine

\page{587}

\note{En tête, au milieu, entre parenthèses et barré : (160). En haut à droite, à l'encre : 289. Le feuillet est monté bas sur sa feuille ; le haut de la vue est la feuille de montage, blanche.}

d'Euler, la courbure n'influe pas sur les manières de vibrer dont
une lame élastique est susceptible et que ce n'est qu'à la circonstance
étrangère dela variation de l'élasticité, qu'il faut attribuer
les différences que M.r Chladni a remarquées entre ses
expériences et la théorie de cet illustre géomètre.

Les mêmes expériences montrent aussi qu'Euler n'auroit
pas dû déroger à son assertion générale sur la non
influence de la courbure, et que c'est à tort qu'il a
exclu le cas des extrémités \struck{appuyés} fixés, de ceux
qui peuvent appartenir à l'anneau.

Je n'ai pas vû de manière de vibrer que je
puisse rapporter dans l'anneau au cas des extrémités
libres, mais ce n'est pas là une raison suffisante de le
rejetter ; car même en supposant qu'il refuse toujours
de se montrer, il suffiroit pour l'exclure que les
figures appartenant à d'autres cas, se formassent plus
facilement, sans que pourtant le cas dont il s'agit, fut
analytiquement impossible.

Lorsque l'on \struck{néglige} n'a pas égard à la variation du \struck{(V. N.os 22 et 23)}
rayon \struck{la plaque} (V. N.os 22 et 23), la plaque circulaire
peut être considérée comme un simple anneau dont la
largeur est mesurée par le rayon de la même plaque et
\note{Le renvoi « (V. N.os 22 et 23) » est d'abord écrit au-dessus de « variation du », puis barré et récrit sur la ligne après « rayon ». « n'a » suit « néglige » barré sur la ligne. Le bas du feuillet est blanc sous la dernière ligne.}

\page{588}

\note{En tête, au milieu, entre parenthèses : (161), non barré. Aucun nombre en haut à droite : c'est le verso écrit du feuillet 289. Le feuillet est monté bas sur sa feuille.}

parceque les sons que peuvent rendre \struck{les} l'anneau\struck{x} \struck{sont}
indépendans de \struck{leur} la largeur\marginal{Note. $\times$ bien entendu qu'il ne s'agit ici que des cas où l'anneau vibre comme une simple lame, car il est susceptible de mouvemens différens dans lesquels la largeur ne peut être négligée}, la plaque doit vibrer alors
de la même manière que si elle étoit véritablement
réduite aux proportions d'un simple anneau.
\note{La note marginale est écrite dans la marge de gauche, en face des lignes 2 à 8, d'une plume plus fine ; elle est appelée par une croix placée au-dessus de « largeur » et répétée en tête de la note. L'x final de « anneaux » est barré. Le premier mot, « les », est repassé.}

Pour vérifier cette théorie, j'ai fait tailler plusieurs
lames circulaires dont le diamètre étoit toujours $=83$ millimetres,
mais dont la largeur étoit, pour les unes, 20 millimetres,
\struck{et} pour les autres \uncertain{35} millimetres. J'ai aussi fait
tailler des plaques \struck{de} circulaires du même diamètre
que les anneaux, voici ceque j'ai observé :
\note{« 35 » : le second chiffre est écrit par-dessus un autre et reste douteux ; le 8 de « 83 » a la forme d'un « \& ». « étoit » (troisième ligne) a sa finale repassée.}

Les lames larges de 20 millimetres ne m'ont pas
donné le cas que j'avois obtenu avec les lames plus
étroites et qui appartient aux extrémités fixés.
\struck{Mais} au reste elles ont donné aisément les deux
\struck{cas} \add{figures} qui sur les lames étroites se rapportent au cas
des extrémités appuyés. L'une d'entre elles a même
donné le cas de deux nœuds, qui répond, comme
on sait, à celui d'un seul nœud sur la lame
droite. \struck{Cette} Ce dernier cas n'a pas été observé
par M.r Chladni sur l'anneau métallique, et
je ne l'ai moi même obtenu que d'une seule
des lames \struck{de ve} circulaires, \struck{de verre}.
\note{« nœuds », « nœud » : l'œ est ouvert et se lit « naud » (voir hand.md). « donné » (« elles ont donné ») porte une lettre finale barrée.}

\page{589}

\note{En tête, au milieu, entre parenthèses et barré : (162). En haut à droite, à l'encre : 290.}

Le son qui accompagnoit la formation des deux nœuds étoit
près du sol plus bas de 2 octaves que l'octave 1.
le son de la figure qui, sur la même lame, \struck{étoit} \add{se} composoit
de deux nœuds, étoit $\text{sol}^{\sharp}$~1 enfin le son qui répondoit
à la figure formée de six nœuds étoit si~3.
\note{Un petit signe, comme un 3 ou un z, est écrit au-dessus de la fin de « étoit », à la première ligne ; il n'est appelé nulle part. « composoit » : la finale est écrite par-dessus une autre. Le dièse de « sol » est écrit en exposant. Le 3 de « si 3 » est de la même forme que celui qu'on a lu, avec doute, dans « $\text{fa}^{\sharp}$ 3 » à la vue 585.}

On voit que, si les sons ne sont pas rigoureusement
conformes à la théorie ; ils \add{ne} s'en éloignent pourtant
pas d'une quantité \struck{supérieure} plus grande que celle
qui peut être attribuée aux causes physiques inappréciables,
et qu'il est permis \struck{de regarder}, parconséquent de regarder
les mouvemens de la lame circulaire large de 20 millimetres
comme parfaitement semblables à ceux d'une lame
circulaire \struck{de même} plus étroite.

Par un hazard heureux, j'ai obtenu absolument les
mêmes sons \add{et les mêmes figures} d'une plaque circulaire et d'un anneau
du même diamètre que la plaque qui ne différoit
d'elle, que par un espace circulaire vide de
11 \struck{\ill{}} millimetres de diamètre\struck{x}. L'une des figures
\struck{l'une} étoit composée de quatre nœuds ou lignes diamétrales
l'autre montroit six semblables nœuds ou lignes
le son étoit $\text{ut}^{\sharp}$~2 dans le premier cas ; dans le
second, le son étoit mi.~3
\note{Après « 11 », un signe barré d'un trait horizontal, peut-être un chiffre (un 7 ?) ; il n'est pas lu. L'x final de « diamètrex » est barré. Le bas du feuillet est blanc.}

\page{590}

\note{En tête, au milieu, entre parenthèses : (163), non barré. Aucun nombre en haut à droite : c'est le verso écrit du feuillet 290.}

Les sons présentent, comme on voit, un intervalle \struck{trop} plus
grand d'un demi ton que ne l'indique la théorie ;
mais dans d'autres expériences l'intervalle s'est trouvé
\struck{juste} \struck{exacte} \add{juste}, et quelquefois aussi trop grand d'un ton
ainsi il doit être considéré \struck{ici} comme suffisamment
exact. \struck{et il est}
\note{« grand » : le g est écrit par-dessus une autre lettre. Au début de la quatrième ligne, « juste » est barré sur la ligne ; au-dessus, « exacte » est barré à son tour et suivi de « juste », non barré.}

On ne peut donc douter que l'équation
$z=\ldots A\,\text{sin}(\beta\varphi+B)$ qui est celle de l'anneau \struck{ne}
\struck{convienne} qui équivaut à une lame dont les extrémités
sont appuyés, ne conviennent également et sous les
mêmes conditions à la plaque circulaire.
\note{« $z=\ldots$ » : un facteur laissé en blanc, comme ailleurs dans le volume. « équivaut » : le mot se lit « equivalt », suivi de quelques lettres barrées. « conviennent » est au pluriel sur la page.}

Il ya plus, \struck{et} comme rien n'empêche d'employer
les arcs de grand cercle qui mesurent les deux courbures
d'une cloche hémisphérique à la place des
coordonnés $x$, et, $y$ et que par une transformation
semblable à celle qui a été pratiquée pour le cas
dela plaque circulaire, on auroit deux facteurs
dans la valeur de \emph{z}, l'un qui contiendroit l'arc
qui remplaceroit le rayon, l'autre qui \struck{conduiroit}
ne \struck{contiendroit que $\text{sin}(\varphi\beta+B)$} seroit autre chose
que $\text{sin}(\varphi\beta+B)$, il sera permis, de négliger aussi
\note{« \emph{z} » est souligné sur la page. Le produit s'écrit ici $\varphi\beta$, là où l'équation de l'anneau, plus haut, porte $\beta\varphi$ ; on garde les deux ordres tels qu'ils sont écrits. Le bas du feuillet est blanc sous la dernière ligne.}

\page{591}

\note{En tête, au milieu, entre parenthèses et barré : (164). En haut à droite, à l'encre : 291. Le feuillet est monté bas sur sa feuille ; le haut de la vue est la feuille de montage, blanche.}

la variabilité du premier facteur et d'appliquer parconséquent
l'équation $z=\ldots A\,\text{sin}(\beta\varphi+B)$ à ce corps
sonore.

Cette \struck{ve} équation $z=\ldots A\,\text{sin}(\beta\varphi+B)$ qui n'est autre
chose que celle de la corde tendue peut donc appartenir,
dans certains cas, à la lame droite, à l'anneau, \struck{à} \add{au tambour circulaire,}
à la plaque circulaire et enfin à la cloche
hémisphérique.
\note{« au tambour circulaire, » est écrit au-dessous de la ligne, sous « à l'anneau », d'une plume plus fine ; le mot barré en fin de ligne est court, peut-être « à ».}

\struck{Il s'en suit que} Dans ceux de ces corps qui sont
élastiques, les sons doivent \add{donc} être proportionnels à
$\beta^2$, c'est-à-dire \struck{à la suite des} \add{aux} quarrés des nombres
dela suite naturelle.

M.r Chladni trouve en effet \struck{que} (V. table p~\uncertain{185})
que, dans le cas où il n'y a sur les plaques
circulaires que des lignes diamétrales, les sons
répondent aux quarrés des nombres 2, 3, 4, 5, 6, \&c.
Il trouve aussi que la cloche d'harmonica\struck{x}
donne, dans les cas analogues, des sons également
proportionnels à la suite des mêmes nombres
(V. table p. 236)
\note{« 185 » : le dernier chiffre, à queue descendante, peut se lire 5 ou 9. Les deux renvois visent des tables, sans doute celles du traité d'acoustique de Chladni cité plus haut ; c'est une inférence. L'x final de « harmonicax » est barré. Le bas du feuillet est blanc.}

\page{592}

\note{En tête, au milieu, entre parenthèses : (165), le 6 repassé ; non barré. Aucun nombre en haut à droite : c'est le verso écrit du feuillet 291. Le feuillet est monté bas sur sa feuille. Cette page est la dernière du mémoire : le texte s'arrête au premier tiers du feuillet, et le reste n'est que l'écriture de la vue 591 vue par transparence.}

Je termine enfin ce long mémoire dans lequel je n'ai pourtant
\struck{encore} \struck{qu'autant} \struck{que commencé} examiné \add{qu'une partie} des différens cas que l'expérience\struck{x}
présente et des cas plus nombreux encore que la
théorie montre possibles. Je m'estimerai heureux
si guidé par l'excellent travail de M.r Chladni
et éclairé par la manière dont Euler a traité un
des cas des corps sonores élastiques, \struck{la classe juge
que je ne me suis pas tout à fait égaré mépris
sur la route qu'il falloit tenir} dans \struck{la carrière
ouverte par cet illustre auteur} \struck{géomètre.}
je suis parvenu à remplir au gré de mes
juges les conditions du programme proposé
par la classe.
\note{Au début de la deuxième ligne, trois leçons barrées se superposent : « encore » et « qu'autant » sur la ligne, « que commencé » au-dessus ; « qu'une partie » est ajouté au-dessus de « examiné », dont l'initiale est repassée. Les trois lignes et demie qui suivent « élastiques, » sont barrées d'un trait continu ; seul « dans » semble épargné par le trait, et la phrase ne le reprend pas. Au-dessus de « illustre », « auteur » est barré et « géomètre » écrit après lui, barré avec le reste. « heureux » est au masculin sur la page.}

\page{593}

\note{Un nouveau morceau commence ici : une mise au net de la première page du mémoire sur les surfaces élastiques dont le brouillon s'ouvre à la vue 398. Aucun numéro de page en tête. En haut à droite, à l'encre : 292. L'écriture est plus régulière et plus grande que dans les pages précédentes, le titre d'une main calligraphiée. La vue 594 photographie une seconde fois ce même feuillet (même texte, mêmes hachures, même nombre 292 ; seul le cadrage diffère) ; elle n'est pas transcrite.}

Mémoire sur la question proposée par la première
classe de l'institut

« Donner la théorie mathématique des vibrations des surfaces
« élastiques et la comparer à l'expérience.

\begin{quote}
Sed longe maximum progressibus Scientiarum, et novis pensis ac
provinciis in iisdem suscipiendis, obstaculum deprehenditur in desperatione
hominum et suppositione impossibilis. \emph{Baco Novum organum, lib. 1,
n. XCII. tom. 1 p. 298, edit. in f.o London 1740.}
\end{quote}
\note{La référence est soulignée ; elle est rendue en italique. Sous l'épigraphe, un trait horizontal au milieu de la page.}

Si on vouloit examiner d'abord la question dans toute sa généralité,
on tomberoit dans des calculs extrêmement compliqués ; mais en se bornant à
considérer le mouvement des surfaces planes dans le cas où il est régulier et
comparable à celui d'un pendule simple, on parvient aisément à une équation
différentielle dont les intégrales particulières sont facilement applicables. J'ose
espérer que, malgré les simplifications que, d'après cette remarque, j'ai crû
pouvoir introduire dans mes formules la classe trouvera la théorie que je
soumets à son jugement, aussi générale que le demande l'état de la
question, car, d'une part, les expériences de M.r Chladni sont toutes comprises
dans le cas des mouvemens réguliers, les seuls en effet qui appartiennent à
l'acoustique, puisqu'eux seuls présentent des lois déterminables ; \& de l'autre
les simplifications auxquelles la supposition de \emph{z} très petit donne lieu, peuvent aisément
se transporter au petit nombre de mouvemens que cet habile physicien a observé
dans les surfaces courbes ; car, comme je le dirai dans la suite,

\struck{cet habile physicien n'a fait qu'un petit \uncertain{nombre} \uncertain{d'essais} sur le son
des surfaces courbes ; et d'ailleurs les géomètres sentiront aisément
que les simplifications auxquelles donne lieu la supposition \emph{z} très petit}
\note{Les trois dernières lignes du texte gardé, depuis « les simplifications auxquelles », sont écrites en alternance avec trois lignes plus anciennes, hachurées de petites croix lettre par lettre : on les donne ici à la suite, en un seul passage barré. Ces lignes hachurées reproduisent la leçon que le brouillon de la vue 398 barrait déjà (« de l'autre, cet habile physicien n'a fait qu'un petit nombre d'essais sur le son des surfaces courbes ; et d'ailleurs les géomètres sentiront aisément que ») ; « nombre » et « d'essais » ne se lisent ici que sous les hachures, d'où le doute. La lettre $z$ est soulignée, deux fois. La phrase s'arrête sur « dans la suite, » au bas du feuillet ; la suite est au revers du feuillet, vue 595.}

\page{595}

\note{Revers écrit du feuillet 292 (vues 593–594) : en tête, au milieu, entre parenthèses : (2), non barré ; aucun nombre en haut à droite. Le texte continue la phrase de la vue 593. Le haut de la page alterne, comme le bas de la vue 593, lignes anciennes barrées et lignes nouvelles écrites entre elles.}

\struck{pourroient se transporter aux surfaces d'une courbure quelconque en}
ces mouvemens se trouvent compris dans une équation fort simple qui convient
\struck{substituant aux différentielles des deux ordonnées rectangulaires \emph{x} et \emph{y}}
à plusieurs genres de corps sonores, et dans laquelle \emph{z} exprime,
\struck{les différentielles des deux courbes qui mesurent les courbures principales
de ces surfaces ; desorte que les valeurs de \emph{z} exprimeroient encore},
comme dans le cas du plan, la quantité dont chaque point des
mêmes surfaces s'éleveroit pendant leur mouvement au dessus de \struck{leur}
situation naturelle.
\note{Les lignes barrées d'un trait sont les anciennes, les autres ont été écrites dans l'interligne ; on les donne dans l'ordre de la page. Ces lignes barrées continuent les lignes hachurées du bas de la vue 593. Le mot barré après « au dessus de » est « leur », avec des lettres récrites par-dessus, peut-être « sa » ; aucun autre possessif n'est écrit, et on ne le supplée pas.}

§. \emph{Recherche de l'équation différentielle dela surface élastique vibrante}
\note{Le titre est souligné sur toute sa longueur ; il est rendu en italique. Le signe qui le précède a la forme d'un 8 suivi d'un point ; on le rend par §, comme le brouillon (vue 399) porte « § 1 ». Aucun chiffre ne le suit ici.}

N.o 1 Pour trouver cette équation, il faut d'abord déterminer quelle doit
être la somme des momens des forces d'élasticité qui agissent dans
toute l'étendue dela surface. Pour cela, je prens quatre points sur
cette surface ; j'imagine deux plans, l'un passant par le premier, le
second et le quatrième de ces points, l'autre passant par le premier, le
troisième et le quatrième des mêmes points ; je nomme $\underline{e}$ l'angle extérieur
formé par un de ces plans et le prolongement de l'autre ; je repre-
sente par $\underline{\beta}$ la petite ligne d'intersection des deux plans et par
$E$ la force qui agit suivant $\underline{\beta}$ en s'opposant à l'inflexion de
la surface et en tendant parconséquent à diminuer l'angle $\underline{e}$.

Ainsi je trouve que la somme des momens est exprimée par $SSE\beta de$ ;
et il ne s'agit plus, pour avoir sa valeur, que de connoître
celles de $de$, $\underline{\beta}$, et $E$.
\note{« $SSE\beta de$ » : le double signe d'intégration est écrit comme deux S, gardés tels quels (voir hand.md). Un petit trait sépare $E$ de $\beta$ ; on ne le rend pas.}

Je commence par la recherche de la valeur de $\underline{e}$.

On sait qu'en représentant par $\underline{C}$ l'ordonnée verticale d'un plan à
son origine, par $A$ et $B$ les tangentes des angles que les traces du
plan dont il s'agit font avec l'horizon, et par $C'$, $A'$ et $B'$ les
quantités analogues relatives à un second plan, le cosinus de l'angle
\note{« traces du plan » : le brouillon (vue 399) ajoutait « sur les plans coordonnés verticaux » ; la mise au net ne l'écrit pas. La page s'arrête sur « de l'angle », au milieu de la phrase ; le bas du feuillet est blanc. La suite n'est pas dans ce lot : la vue 596 porte une page numérotée (23), de la même main.}

\page{596}

\note{En tête, au milieu, entre parenthèses : (23), non barré. En haut à droite, à l'encre : 293. Même main régulière qu'aux vues 593 et 595. Le feuillet est monté bas sur sa feuille ; le haut de la vue est la feuille de montage. La page commence au milieu d'une phrase, dont le début n'est pas dans ce lot : la page précédente de cette série, (2), s'arrête à la vue 595 sur « de l'angle ».}

équations, j'examinerai les conséquences que l'on peut en tirer dans les
hypothèses les plus simples. Je supposerai donc d'abord que la surface
vibrante soit rectangulaire, que le plan de cette surface en repos soit
pris pour celui des $\underline{x}$ et $y$ afin que l'origine soit à un des points
de repos de la même surface, enfin que chacun des côtés soit parallèle
à un des axes des $\underline{x}$ et $\underline{y}$. Cela posé, la suite des points qui servent
de limites aux valeurs de $\underline{x}$, c'est-à-dire ceux qui appartiennent aux
côtés parallèles à l'axe des $y$ devront satisfaire aux conditions
\[ S\frac{1}{dy}\left\{\frac{ddz}{dx^2}+\frac{ddz}{dy^2}\right\}\delta\!\left(\frac{dz}{dx}\right)=0, \quad S\frac{1}{dy}\,\frac{d\left\{\frac{ddz}{dx^2}+\frac{ddz}{dy^2}\right\}}{dx}\,.\,\delta x=0, \]
en vertu de l'un ou l'autre des quatre couples d'équations :
\[ z=0 \text{ et } \frac{dz}{dx}=0\,; \quad z=0 \text{ et } \frac{ddz}{dx^2}+\frac{ddz}{dy^2}=0\,; \]
\[ \frac{d\left\{\frac{ddz}{dx^2}+\frac{ddz}{dy^2}\right\}}{dx}=0 \text{ et } \frac{dz}{dx}=0\,; \quad \frac{d\left\{\frac{ddz}{dx^2}+\frac{ddz}{dy^2}\right\}}{dx}=0 \text{ et } \frac{ddz}{dx^2}+\frac{ddz}{dy^2}=0. \]
\note{« en repos » est écrit d'une encre plus lourde, par-dessus d'autres lettres, et « soit » est traversé d'un trait qui peut être un trait de liaison ; on ne le donne pas comme biffé, la phrase en ayant besoin. Les signes $S$ sont gardés tels qu'ils sont écrits, et le facteur est bien $\frac{1}{dy}$ sur la page ; le $\delta$ est la lettre bouclée de la variation. La fin de la première condition porte « $.\,\delta x$ » ; la lettre après $\delta$ est lue $x$.}

L'analogie est frappante entre ce que l'on trouve ici et cequi
a lieu pour les extrémités dela simple lame. On voit en effet
que les deux derniers couples ne peuvent appartenir qu'aux points
libres du contour, tandis que les deux premiers conviennent aux points
de repos du même contour. Cependant on ne peut plus regarder
le point auquel appartiendroient les conditions $z=0$ et $\frac{dz}{dx}=0$
comme fixé parcequ'elles ne suffisent pas pour que le plan tangent
tombe dans le plan de $\underline{x}$ et $y$ et qu'il faudroit pour cela
que l'on eut encore $\frac{dz}{dy}=0$.

On voit de même que la suite des points qui servent de
limites aux valeurs de $y$, c'est-à-dire ceux qui appartiennent
aux côtés parallèles à l'axe des $\underline{x}$ doivent satisfaire aux conditions,
\[ S\frac{1}{dx}\left\{\frac{ddz}{dx^2}+\frac{ddz}{dy^2}\right\}\delta\!\left(\frac{dz}{dy}\right)=0, \quad S\frac{1}{dx}\,\frac{d\left\{\frac{ddz}{dx^2}+\frac{ddz}{dy^2}\right\}}{dy}\,\delta z, \]
en vertu de l'un ou
\note{La seconde condition s'arrête sur « $\delta z$, » sans « $=0$ » ; elle est donnée telle qu'elle est écrite. Là où la première série portait $\delta x$, celle-ci porte $\delta z$. La page s'arrête sur « en vertu de l'un ou » ; la suite est au revers, vue 597.}

\page{597}

\note{Revers écrit du feuillet 293 (vue 596) : en tête, au milieu, entre parenthèses : (24), non barré ; aucun nombre en haut à droite. Le feuillet est monté bas sur sa feuille. Le texte continue la vue 596.}

l'autre des quatre couples d'équations : $z=0$ et $\frac{dz}{dy}=0$ ; $z=0$ et $\frac{ddz}{dx^2}+\frac{ddz}{dy^2}=0$ ;
\[ \frac{d\left\{\frac{ddz}{dx^2}+\frac{ddz}{dy^2}\right\}}{dy}=0 \text{ et } \frac{dz}{dy}=0\,; \quad \frac{d\left\{\frac{ddz}{dx^2}+\frac{ddz}{dy^2}\right\}}{dy}=0 \text{ et } \frac{ddz}{dx^2}+\frac{ddz}{dy^2}=0. \]

Puisqu'il s'agit ici des surfaces rectangulaires, on voit que les points
situés aux quatre angles sont assujettis à quatre conditions : car ils
appartiennent également aux côtés parallèles à l'axe des $\underline{x}$ et aux
côtés parallèles à l'axe des $y$. Je nommerai points de \emph{double
limite} tout point qui remplira à la fois les conditions des limites
des valeurs de $\underline{x}$ et celles des valeurs de $y$. Ainsi les équations
$z=0$ et $\frac{ddz}{dx^2}+\frac{ddz}{dy^2}=0$ qui forment le caractère analogue à
celui qu'Euler attribue aux points d'appui, ne pourront appartenir
qu'aux points de \emph{doubles limites}. On aura encore un point de
\emph{double limite}, si on peut satisfaire à la fois aux deux couples
d'équations $z=0$ et $\frac{dz}{dx}=0$ ; $z=0$ et $\frac{dz}{dy}=0$, et le plan
tangent à ce point tombera dans le plan des $\underline{x}$ et $y$. Au contraire
le plan tangent sera seulement parallèle au plan des $\underline{x}$ et $y$
lorsque les conditions du point de \emph{double limite} seront les deux
couples d'équations $\frac{d\left\{\frac{ddz}{dx^2}+\frac{ddz}{dy^2}\right\}}{dy}=0$ et $\frac{dz}{dy}$ ;
\[ \frac{d\left\{\frac{ddz}{dx^2}+\frac{ddz}{dy^2}\right\}}{dx}=0 \text{ et } \frac{dz}{dx}=0. \]
\note{« double limite » est souligné chaque fois ; il est rendu en italique. « sont assujettis » : « sont » est écrit par-dessus un autre mot. « Euler » est écrit de son E en forme de L. Dans le premier couple de la dernière formule, « $\frac{dz}{dy}$ » est écrit sans « $=0$ » ; on le donne tel quel.}

Il est clair que tout point qui satisfera aux conditions des
limites appartiendra à une ligne de limite analytique : mais de plus
on peut observer que si une pareille ligne de limite satisfait, par exemple,
aux conditions $S\frac{1}{dy}\left\{\frac{ddz}{dx^2}+\frac{ddz}{dy^2}\right\}\delta\!\left(\frac{dz}{dx}\right)=0$, $S\frac{1}{dy}\,\frac{d\left\{\frac{ddz}{dx^2}+\frac{ddz}{dy^2}\right\}}{dx}\,.\,\delta z=0$.
et que pour un ou plusieurs de ses points on ait $dy=0$,
\note{La seconde condition porte ici $\delta z$, là où la même condition, à la vue 596, porte $\delta x$ ; chacune est donnée telle qu'elle est écrite. Le « 0 » de « $dy=0$ » est un peu fermé et pourrait être lu « a » ; le sens donne 0. La page s'arrête sur cette virgule, au bas du feuillet ; la suite n'est pas dans ce lot.}

\page{598}

\note{Un autre état du même début : un brouillon de la première page du mémoire sur les surfaces élastiques, sans numéro de page en tête, très corrigé. En haut à droite, à l'encre : 294. Le texte s'arrête aux deux tiers du feuillet ; le bas est blanc.}

Mémoire sur cette question proposée par la première classe de
l'institut.

« Donner la théorie mathématique des vibrations des surfaces élastiques
« et la comparer à l'expérience ?.

Si on vouloit embrasser d'abord cette question dans toute sa généralité
on tomberoit dans des calculs extrêmement compliqués, mais en se bornant
d'abord \struck{aux cas des surfaces planes, et en ne s'occupant que de ceux
des mouvemens de semblables surfaces qui peuvent être comparés à
celui d'un pendule simple} à considérer le mouvement des \struck{seules}
surfaces planes dans le cas où il est régulier et comparable
à celui d'un pendule simple on parvient aisément, comme je vais
le démontrer, à une équation \add{différentielle} \struck{du quatrième degré} \uncertain{fort} simple.
\struck{à la vérité cette équation n'est pas intégrable en}
J'ose espérer que la classe \struck{jugera} \add{trouvera} la théorie que j'ai l'honneur de soumettre à son jugement
\marginal{aussi générale que le \struck{veut} \add{demande} l'état dela question car d'une part}
\note{« seules » est écrit une première fois au-dessus de « aux cas des », dans la partie barrée, et une seconde fois en fin de ligne après « le mouvement des », où il est barré. « fort » est taché d'un pâté et peut être biffé. La ligne « J'ose espérer … à son jugement » est serrée dans l'interligne, au-dessus de la ligne barrée « à la vérité … intégrable en » ; « trouvera » est écrit au-dessus de « jugera », barré. La suite de la phrase est dans la marge de gauche, en quatre lignes, sans signe de renvoi ; « demande » est écrit au-dessus de « veut ».}

Et On \struck{Au reste} \add{sait que} les expériences de M.r Chladni, sont \add{\struck{, comme on sait}} \struck{toutes ou presque toutes}
comprises dans le cas \struck{puisqu'il ne s'occupe jamais que des mouvemens}
des mouvemens réguliers \add{les seuls en effet qui appartiennent à l'acoustique c'est à dire \struck{d'un} celui où des} \struck{puisqu'il ne s'occupe que de ceux qui donnent
lieu à des} sons perceptibles \struck{et à des} \add{accompagnent des} figures \struck{régulières} constantes.
\note{« Et On » : les deux mots sont écrits d'une encre plus lourde en tête de ligne, « On » repassé ; « sait que » est au-dessus de « Au reste », barré. Les deux additions interlinéaires « les seuls en effet qui appartiennent à l'acoustique » et « c'est à dire d'un celui où des » sont écrites l'une au-dessus de l'autre, au-dessus de « réguliers puisqu'il » ; « accompagnent » est au-dessus de « et à des », barré, en fin de ligne ; on les place là où la phrase les demande. Le « des » qui précède « figures » est lu dans « et à des », qui n'est barré que sur ses deux premiers mots : c'est une lecture.}

\add{\struck{et} de l'autre \uncertain{habile}} d'ailleurs \struck{cet} physicien n'a fait qu'un petit nombre d'essais sur le
son des surfaces courbes \struck{ainsi} \struck{desorte que le cas auquel s'applique
la théorie que j'ai l'honneur de soumettre au jugement de la classe
a je crois toute la généralité convenable à l'état de la question.}
\note{Au-dessus de « d'ailleurs cet », une addition : « et de l'autre », « et » barré, suivi d'un mot qui se lit « habile » et chevauche « l'autre ». « ainsi », barré, est au-dessus de « desorte », barré. Les deux lignes suivantes, « la théorie … question. », sont hachurées de petites croix lettre par lettre ; « je crois » y est lu sous les hachures.}

\add{et d'ailleurs} \struck{ma} \struck{Au reste} les géomètres sentiront aisément que la simplification à la-
quelle donne lieu la supposition de $\underline{z}$ très petit pourroit se transporter
aux surfaces d'une courbure quelconque en substituant aux.
\note{Le premier mot barré de la ligne est court, peut-être « ma » ; « Au reste » le suit, barré aussi ; « et d'ailleurs » est écrit au-dessus. La page s'arrête sur « aux. », le point écrit ; la phrase n'est pas achevée.}

\page{600}

\note{En tête, au milieu, entre parenthèses et barré : un nombre commençant par 6, lu (61) ; le second chiffre est pris dans un pâté, mais la vue suivante, hors de ce lot, porte (62) en tête, ce qui appuie la lecture. En haut à droite, à l'encre : 295. D'une autre suite que les vues 593–598 : la page commence au milieu d'un raisonnement sur les intégrales de la plaque et sur les figures de M.r Chladni, et son début n'est pas dans ce lot. Le feuillet est monté bas sur sa feuille.}

En effet j'ai vu bien des fois les lignes dues aux facteurs périodique, se
séparer et se courber : mais M.r Chladni s'est bien gardé de tenir
compte de cet accident ; et malgré l'ignorance où il étoit de la
théorie, il a deviné parmi les distorsions qui sembloient être de même
nature qu'elles étoient celles que l'on ne devoit attribuer qu'à
quelqu'imperfection de la plaque vibrante.

A l'égard de l'intégrale (++), il est aisé de voir que, quoique
les figures qui pourroient y être rapportées, doivent présenter des lignes
droites nodales dans le sens des deux axes, les sons qui accompagneroient
ces fig., devront être proportionnels à des nombres qui suivent entr'eux
la même loi que ceux qui conviennent aux différens cas du mouvement
de la simple lame dont les extrémités sont appuyés ou de celle
dont les extrémités seroient libres en vertu des conditions $\frac{dz}{dx}=0$ et $\frac{d^3z}{dx^3}=0$.

Il est clair que les séries qui entrent dans l'intégrale (++), ne peuvent
jamais se terminer ; ainsi elle ne peut \struck{être d'une application} facile :
c'est pourquoi je ne m'y arreterai pas.
\note{L'étiquette de l'intégrale est faite de deux petites croix entre parenthèses ; on la rend (++). Un long trait oblique, d'une plume fine, descend de la ligne « A l'égard de l'intégrale (++) » jusqu'au-dessous de « je ne m'y arreterai pas » et traverse ces deux alinéas : il semble les annuler, et on le dit ici plutôt que de les donner comme barrés. « être d'une application » est barré sans que rien le remplace. « accompagneroient » a sa finale barrée d'un trait.}

L'intégrale (\uncertain{++}) qui appartient aux surfaces qui ne sont périodiques
que dans le sens des $y$, \add{renferme une série qui ne peut jamais se terminer} \struck{présente la même singularité que l'intégrale
(++), c'est-à-dire que les sons qui accompagnent les figures qui s'y
rapportent, sont également proportionnels aux nombres qui conviennent
aux mouvemens de la simple lame appuyée à ses deux extrémités.
La série qui entre dans cette intégrale,} ne peut non plus jamais se
\struck{terminer} ; \struck{et} il est parconséquent impossible de déterminer les
figures qu'elle peut expliquer, sans avoir recours au calcul.
\struck{et je ne m'y arreterai pas}
\note{L'étiquette de cette intégrale a ses deux croix réunies par un second trait, comme une petite grille ; on ne sait si c'est la même que la précédente. « renferme une série qui ne peut jamais se terminer » est écrit au-dessus de « présente la même singularité », d'une plume plus fine. Le trait qui barre les lignes suivantes s'interrompt sur « ne peut non plus jamais se » ; on donne la page telle qu'elle se voit, bien que la phrase gardée n'appelle pas ces mots.}

Cependant je crois pouvoir assurer que les fig. \uncertain{796} et 86 de M.r
Chladni doivent être rapportées à cette intégrale, dans les suppositions
successives, $m=5$, $m=6$.
\marginal{\add{A} l'égard des sons qui accompagneroient ces figures \add{on} voit qu'ils peuvent être \add{les} mêmes que ceux qui \add{sont} donnés par les intégrales périodiques dans le sens des $\underline{x}$ et dans celui des $y$ car les nombres \struck{$(m^2+n^2)$}, $(m^2+n^2)(\Pi-1)(\Pi+\Pi'-2)$ \struck{\ill{}} se réduisent à la \struck{\ill{}, être deux même} (en prenant pour $\Pi$ et $\Pi'$ des valeurs convenables) \struck{et de la} forme}
\note{« 796 » : le dernier chiffre a exactement la forme du 6 de « 86 » qui suit ; un numéro de figure aussi élevé surprend, et ce peut être un 79 suivi d'une lettre ; d'où le doute. La vue 601, hors de ce lot, écrit le même numéro de la même façon. La note est écrite dans la marge de gauche, au bas du feuillet, en six lignes, puis continue sous la dernière ligne du texte, en travers de la page, jusqu'au bord ; plusieurs mots en sont perdus dans la reliure, les mots donnés comme additions sont suppléés par le sens. « successives ; $m=5$, $m=6$. » est écrit sur la même ligne que la fin de la note, dont il est séparé par « se réduisent à la » ; on le rattache au texte par le sens. Les lettres $\Pi$, $\Pi'$ sont le grand $\pi$ à jambes hautes que les lots précédents rendent ainsi (voir hand.md) ; ici la lettre ressemble à « rr », et on l'offre sous réserve. Au-dessous, barrés, un mot court puis « pendant », et « être deux même » ; le mot barré avant « se réduisent » ne se lit pas. L'ordre des derniers mots de la note (« se réduisent à la », la parenthèse, « forme ») est celui que la page suggère, non un ordre assuré. La page s'arrête sur « forme », en bas du feuillet ; la suite n'est pas dans ce lot.}

\end{document}
